\chapter*{前\hspace{1em}言}
\markboth{前\hspace{1em}言}{前\hspace{1em}言}

2016 年夏天，我在小学机器人兴趣班第一次看到自己写的 C 代码让小车动起来。屏幕上一行行冰冷的文字，下载到单片机里竟然变成了轮子的转动。那一刻的兴奋，我到现在还记得。

从那以后，我就一直在这条路上走。我的专业是计算机科学与技术，Arduino、STC51、STM32、K230 是我一路摸过来的路标。从学生到工程师，这本书，就是这一路踩过的坑、烧过的板子、和深夜没睡着的那些晚上攒下来的东西。

起初我没想过要写书。这事儿还得从在实验室带学弟学妹说起——当时为了培训他们搞协程调度器，我吭哧吭哧做了一套PPT，顺手整理了一本薄薄的小培训教程。但在教的过程中我发现，很多底层细节光靠一份课件根本讲不透；加上自己平时学的时候，同一个知识点也总翻来覆去地查：状态机怎么写、FIFO 怎么实现、菜单框架怎么搭……每次重新搞懂都要花不少时间。于是我开始给自己、也给他们记笔记，把“搞懂一次之后不想再忘记”的东西一点点往里补充。笔记越写越多，越写越系统，最后就变成了这本书的雏形。

这些章节的顺序，就是我自己爬过来的路线——从最基础的地方开始，一步一坑地往前走，每章都比前一章难一点，但又不至于让你卡在半路动不了。如果你读一遍就能全记住，那这本书不是写给你的。我写的是那种“读的时候骂骂咧咧，调通之后又觉得值了”的东西。

第17章、第18章和第19章跟前面那些都不一样。它们讲的是时间不够的时候怎么活下来，怎么用 AI 给自己打工，以及怎么把个人经验变成能重复执行的产线。第20章移到了全书最后，讲这本书和我自己的来路，让故事在技术之后收束。如果你正在赶交期、正在被一个调不通的外设折磨，或者正在犹豫要不要继续做嵌入式，这几章可能比单独的代码更能帮到你。

\section*{关于书里的代码}
这本书里的代码，大多数是 C 语言风格的伪代码，用于对照，不是直接参考实现。你可以看结构、看流程、看判断，但不要指望复制粘贴就能跑。书里的代码在 PC 上能跑，是为了让你看懂逻辑；到了单片机上，能不能用、该怎么改、要不要换方案，是你的决策。没有逐一到单片机上跑通，因为时间有限、平台太多。板级验证是读者自己的第一步。书给的是骨架，肉要读者自己长。

\section*{为什么前面要放这些项目}

你可能会问：这本书讲 Agent、讲产线、讲 AI 协作，为什么前面 16 章还要花那么多篇幅讲 C 语言、内存池、状态机、串口、FIFO、菜单、命令解析、PID 和 FFT？

因为知己知彼，才能百战百胜。你得知道这些项目是怎么来的，你才能改它。你得知道 FIFO 为什么牺牲一个槽位，你才能让 AI 帮你改缓冲逻辑。你得知道状态机为什么拆成那样，你才能判断 AI 生成的状态机对不对。你得知道内存池的边界在哪，你才能看出 AI 写的分配器有没有越界。

AI 时代，最重要的不是 AI，是你自己的核心。核心不牢，你连让 AI 写什么都说不清楚。你不知道代码该写在哪，不知道这个功能该加在哪个模块，不知道这段逻辑该放中断还是主循环——那你叫 AI 写，它写出来的东西你也判断不了对错。简单说：你得知道一些常识。你得知道代码大概长什么样，你得知道一个功能通常加在哪里，你得知道改一个地方会影响哪些地方。前 16 章，就是在给你这些常识。有了这些常识，你才能带着 Agent 走 Make it work $\rightarrow$ Make it right $\rightarrow$ Make it fast。没有这些常识，你只能被 Agent 牵着走。

\section*{基础不牢，工具也救不了你}

这本书不会教你 Codex 怎么部署、DeepSeek 怎么注册、Claude 怎么充值。不是我不知道，是我不想写。因为这些是基础中的基础。你得学会使用互联网资源，自己查。官网有文档，社区有帖子，搜索引擎有答案。你连这个都搞不定，那后面的事你更搞不定。

打个比方：你想学 3D 建模，但 3D 建模软件你装不上、界面你看不懂、快捷键你记不住，那我劝你最好别学。不是建模难，是你连第一步都不愿意自己走。AI 工具也一样。注册、配置、API key、环境变量，这些东西网上全是教程。你花半小时能查明白的事，非要等人喂到嘴边，那你就永远只能被别人牵着走。这本书教的是判断，不是操作。操作你自己查，判断我陪你练。

书里的代码风格可能不太像教科书。我喜欢把东西写得很“直”：能不用宏就不用宏，能不封装就不封装，能用 \texttt{if-else} 说清楚的就不用设计模式。这种风格在工程上未必最优，但它有一个好处——每一行的意图都很清楚，不会让你在“这语法什么意思”上卡住。等读懂了直白的写法，再去学优雅的写法，比反过来容易得多。

最后，说说这本书怎么来的。我先根据自己的笔记和项目经验，搭好了全书结构；然后花了一段时间逐章展开，把每个知识点用清晰的文字和可运行的代码写出来。写完之后又反复通读了好几遍，删掉不准的说法，补上只有在真硬件上跑过才知道的细节，调整语言让它读起来顺一些。这个过程里，我既是作者，也是自己的第一个读者——凡是我读着费劲的地方，都重新写过。

这种写法的好处是快——从动笔到初稿完成只花了两三个月；坏处是书的“口感”可能不太均匀，有些段落像深夜赶工的手笔，有些又像反复打磨后的产物。如果你读到了不舒服的地方，大概率是我赶工留下的痕迹，欢迎你提出来，我后面改。

感谢每一位翻开这本书的读者。如果书里的某个思路、某段代码能在某个深夜给你一点启发，或者让你觉得“原来这条路真的有人走过”，那它存在的意义就实现了。

\section*{关于书名}

原来这本书叫《嵌入式工程师的荒野求生笔记》，讲的是一个人、一块板子、一个交期，怎么活下来。后来我发现，活下来之后还有一件事：怎么不再靠一个人硬扛，怎么用 Agent、分支、PR、全量回归，把一次性的求生，变成可复制的产出。所以书名改成了《嵌入式工程师的基础知识与产线构建》。荒野是前 17 章，产线是 18、19 章。前者教你别死，后者教你持续活着。

愿你所有的调试，都亮起绿灯。

\rightline{—— 写于深夜的深圳出租房}
